%% The following is a directive for TeXShop to indicate the main file
%%!TEX root = diss.tex

%=======================================
\chapter{Algebraic Kernel method}
\label{ch:AKM}
%=======================================
%=======================================
\section{Non-interacting case}
%=======================================
%=======================================
\subsection{Walks with small steps}
%=======================================
At the core of the Algebraic Kernel method we have the idea that we can take the master functional equationapply group actions on it from $G$ and obtain non-trivially different equations for each $g\in  G$. After we sum these equations in a nice way we manage to eliminate the $Q_1^0,Q_2^0$ and $\Qf{0,0}$ and their orbits leaving only the orbit of the kernel term $KQ$ and a term on the right-hand side. What makes this possible is pre-multiplying the whole equation by $x_1\ldots x_2$. This also works in higher dimensions as we'll show later. Although the generalization has been done before we will present it in a compact and easy-to-generalize way. Instead of thinking just in terms of equations and summing them we'll encapsulate the process in an operator $D_G$ that will take the functional equation and eliminate the terms we don't want by applying the group actions in a nice way. We start by defining $D_G$ in order to formalize the process in the Algebraic Kernel method.
\begin{definition} Let $G$ be a group represented and generated by $g_k$'s preserving the kernel $K$ for a model for small steps. Then let $D_G$ be the operator defined as
\begin{align}
	D_G = \sum_{g\in G} (-1)^{|g|}g
\end{align}
in the sense of group actions.
\end{definition}
In two dimensions we can obtain the whole group by iterating over the two generators 1, $g_i$, $g_jg_i$, $g_ig_jg_i$. It's better to write something like $g_{iji}$ instead of $g_ig_jg_i$ so we have the following.
\begin{definition} For a sequence $L = (l_1,l_2\ldots)$ of positive integers define 
\begin{align}
	g_L = g_{l_1l_2\ldots} = g_{l_1}g_{l_2}\ldots
\end{align} 
whenever this makes sense - when the list is finite or when there is a convergence in the case of an infinite list.
\end{definition}
In the case of a finite group in two dimensions we have $(g_1g_2)^n = 1$ for some minimal $n$ which gives the following lemma so we can write the operator $D_G$ in a neater format.
\begin{lemma} For a two-dimensional random walk model with small steps possessing a finite group of order $2n$ and generators $g_1$, $g_2$ we can write the operator $D_G$ as
\begin{align}
	D_G = (1+g_{ij}+\cdots+(g_{ij})^{n-1})(1-g_i)
\end{align}
\end{lemma}
\begin{proof} We expang the identity in the statement of the lemma
\begin{align}
	D_G &= (1+g_{ij}+\cdots+(g_{ij})^{n-1})(1-g_i)\\
	        &= 1- g_i +g_{ij}-g_{iji} + \cdots - (g_{ij})^{n-1}g_i
\end{align}
giving us one of the possible orbits
\end{proof}
Here are some examples of small groups
\begin{example} When $|G| = 4$ the elements commute so the orbit is $1, g_1,g_{12} = g_{21}, g_2$ and the operator is
\begin{align}
	D_G = 1 - g_1 - g_2 + g_{12} = (1-g_1)(1-g_2)
\end{align}
When $|G| = 6$ the orbit is $1,g_1,g_{21},g_{121}=g_{212},g_{12},g_2$. We can write the operator as
\begin{align}
	D_G &= (1+g_{12}+g_{1212})(1-g_1) \\
	&=  1+g_{12}+g_{1212} - g_1 - g_{121} - g_{12121} \\
	&= 1 -g_1 - g_2 + g_{12}+g_{21}-g_{121}
\end{align}
\end{example}
We can generalize the lemma for a multi-dimensional group $G$.
\begin{lemma}Let $G$ be a group for a model with small steps. Then we can write the operator $D_G$ in terms of $g_k$ as $D_G = D_H(1-g_k)$ for a subgroup $H$ of $G$.
\end{lemma}
\begin{proof} Simpler proof: each element in the fixed representation of $G$ either ends in $g_k$ or doesn't. We have a bijection between the two sets by using the involution $g_k$ applied on the right. The elements ending in $g_k$ have opposite parity to those not ending on $g_k$ so we can split $G$ into two parts $H$ and $Hg_k$. Splitting the defining sum for $G$ into two parts and keeping track signs via $e_g = (-1)^|g|$ leads to
\begin{align}
	D_G &=\sum_{g \in H} e_g g + \sum_{g \in Hg_k} e_g g\\
	&=\sum_{g \in H} e_g g + \sum_{g \in H} (-1)e_g g g_k\\
	&= \left(\sum_{g \in H} e_g g\right)(1-g_k)\\
	&= D_H(1-g_k)
\end{align}
I'll have to provide a different proof maybe. Take $H = G/\langle g_k \rangle$. Since $g_k$ is of order 2 the resulting group has $|G|/2$ elements. Consider the elements of $G$ that commute with $g_k$, that is $h\in G$ where $g_k h = h g_k$. Multiplying by $g_k^{-1} = g_k$ we see that $h = g_k h g_k$. Let $H$ be the subgroup of $G$ generated by those elements. There are two cosets $H$ and $Hg_k$.
\end{proof}
Thus we have
\begin{lemma} The operator $D_G$ eliminates the term $t S_I^{-\one}Q_I^\oo$ for $I \neq \emptyset$.
\end{lemma}
\begin{proof} For $i \in I$ write $D_G = D_H (1- g_i)$ by the previous lemma and then apply it. Since the variable $x_i$ is not present in the $S_I^{-\one}Q_I^{\oo}$ term we see that $(1-g_i)$ annihilates it and the proof is complete.
\end{proof}
This leads us to the following theorem
\begin{theorem} For a model with small steps we have
\begin{align}
D_G(x_1\ldots x_d KQ) = D_G(\mon xd)
\end{align}
where the RHS may be zero sometimes.
\end{theorem}
\begin{proof}Multiplying through by $\mon xd$ gives us the situation in the lemma and so we obtain the eliminations we wished for.
\end{proof}

\begin{corollary} Suppose that there is a group element $h \in G$ with odd parity that permutes the variables $x_1,\ldots,x_d$ in some way and has order 2. Then the regular Algebraic Kernel method fails. That is
\begin{align}
D_G(\mon xd KQ) = 0
\end{align}
\end{corollary}
\begin{proof} Since $h$ is an involution the same method applies as before and we can factor $D_G = D_H(1-h)$ for some $H \subseteq G$. The action $1-h$ eliminates the monomial $\mon xd$.
\end{proof}

\begin{example} The 2D models where the Algebraic Kernel method fails in its straight-forward construction are precisely those where we have a permutation as in the corollary. That is $h(x_1,x_2)= h(x_2,x_1)$ and $|h| = 3$.
\end{example}

\begin{example} In Steve's paper there is a model with group size 12 for which the Algebraic Kernel method fails. For that model one of the group element is a permutation of two of the variables - in that case $x_2$ and $x_3$. Strangely enough $G$ contain the group seen in the previous example for the variables $x_2$ and $x_3$. This means that every higher-dimensional model that has group $G$ containing the group above will fail in the regular Algebraic Kernel method.
\end{example}
A natural question is to consider taking $D_H$ instead of $D_G = D_H(1-h)$ and checking if we can perfom a half-orbit summation in general. The answer is yes. The operato $D_H$ eliminates some of the $Q$ variables and introduces others but the result is a function symmetric under permutation by $h$. One may hope to perform a method analogous to the half-orbit summation method to prove algebraicity or D-finite-ness

There is another way we can prove that $D_G$ does what it advertises. We start with the kernel form $EKQ = 1$ where we make the substitution $Q = W/(\mon xd K)$ and re-write the equation as
\begin{example}
\begin{align}
	E' W = \mon xd,\quad\quad E_k' = 1 - [x_k^0]
\end{align}
\end{example}
Therefore
\begin{align}
	W = \mon xd -(1-E')W=\mon xd - \sum_{\emptyset \neq I \subseteq [d]}W_I^\oo
\end{align}
We immediately see that
\begin{align}
	D_G(\mon xd KQ) = D_G(W) &= D_G\left(\mon xd - \sum_{\emptyset \neq I \subseteq [d]}W_I^\oo\right) \\&= D_G(\mon xd)
\end{align}
since the $W_I^\oo$ terms have missing variables. Let's think about this for a moment. The form of $W$ is the same for all random walk models with small steps so if there are two models whose operators both annihilate expressions of the form $W_I^\oo$ for all non-empty $I\subseteq [d]$ then both models have equations of the above type with each other's operators! In particular we may take the group $G$ of the highly-symmetric models and apply it via the operator $D_G$ for any other model and we obtain an equation of the above type! This observation suggests looking into this idea further.
\begin{definition} Let $g_k(x_k) = \ol{x}_k$. Define the operator $\Delta_k = 1 - g_k$ for this specific $g_k$. Let $\Delta = \Delta_1\cdots \Delta_d$.
\end{definition}

\begin{theorem} Let $G$ be the group generated by $g_k(x_k) = \ol{x}_k$ for $1 \leq k \leq d$. Let $\Delta_k = 1 - g_k$. Then we have $D_G = \Delta$
and
\begin{align}
	\Delta(\mon xd KQ) = \Delta(\mon xd) = (x_1 - \ol{x}_1)\cdots(x_d-\ol{x}_d)
\end{align}
for any small step random walk model in dimension $d$ with no interactions.
\end{theorem}
\begin{proof} It's enough to see that $\Delta$ annihilates any term of the form $W_I^\oo$ for $I \neq \emptyset$. Since the operatos $\Delta_k$ commute we only need to apply one $\Delta_i$ for some $i$ first and find
\begin{align}
	\Delta_i(W_I^\oo) = W_I^\oo - g_i(W_I^\oo) = 0
\end{align}
since the term is missing variable $x_i$. The rest follows by linearity.
\end{proof}
This may not be of great value for showing a specific model has a D-finite solution because we cannot factor out the $K$ term from the LHS in general. However the idea is important because it allows us to prove the following startling corollary.
\begin{corollary} All random walk models of dimension $d$ satisfy the functional equation $E'(\mon xd KQ) = \mon xd$. Therefore the only difference between the models is in the initial conditions that go with this functional equation.
\end{corollary}

We'll extend these ideas in the next section. But first we'll look into generalizing the idea of the orbit summation operator.
\begin{definition} Let $G$ be a group acting on $\row xd$. Then we define a weighted orbit summation operator to be any operator of the form
\begin{align}
	D = \sum_{g\in G}R_g g
\end{align}
where $R_g$ are functions in the variables $\row xd$.
\end{definition}

The orbit summation operator is an example of this with $R_g' = 1$ for those groups considered so far. This kind of operator captures a great dead of things we may try to do. For example, we may pre-multiply the whole functional equation by some function $T$. Then we can think of this function $T$ as being a part of a weighted orbit summation operator since
\begin{align}
	D(Tf) = \sum_{g\in G} R_g g(T f) = \sum_{g\in G} (R_g g(T)) g(f) =  \sum_{g\in G} R_g' g(f).
\end{align}
Also repeated applications of $D$ are also of the same form since $G$ is a group.
\begin{prop} For a group $G$, the set of weighted orbit summation operators forms a semi-group under the usual operations.
\end{prop}
\begin{proof}We'll just check that the set is closed under composition. We have
\begin{align}
	D_1D_2 = \sum_{g\in G}\sum_{h\in H} R_g gT_h h = \sum_{g\in G}\sum_{h\in H} R_g T_{gh} gh \sum_{g\in G}(\sum_{h\in H} R_{gh_{-1}} T_{g}) g
\end{align}
which is also an operator of this form once we collect terms corresponding to the same group elements.
\end{proof}

We hope to apply a weighted operator $D$ to a problem in order to gain more insight into the structure of the generating function $Q$. Since all models with small steps that have finite groups also have a regular orbital operator we have two main interesting cases
\begin{itemize}
\item[(1)] when $D_G(\mon xd) = 0$
\item[(2)] when $G$ is infinite.
\end{itemize}
The question is whether moving to a weighted operator will lead to something that works. Unfortunately in the first case the answer is no - in general there is no other weighted orbital operator that doesn't also eliminate $\mon xd$. The second case remain unclear.

Here's how the method works. If we apply the operator $D$ on the master functional equation we should only have the kernel term left. That is $D$ needs to annihilate all the $Q$ terms that are missing variables. The multitude of random walk models with small steps is huge so its unwieldy to use the master functional equation. However we have the fundamental functional equation for the function $W = \mon xd KQ$ so it's enough to work with the fundamental functional equation instead. To go back to the master functional equation case we recognize that $\mon xd$ just make it a weighted operator and we can then de-weigh it by pre-dividing by $\mon xd$. The resulting operator is also weighted.

\begin{lemma}A weighted operator $D$ exists such that $D( KQ) = D(1)$ if and only if a weighted operator $D'$ exists such that $D'(f_I^\oo)$ for generic functions $f_I^\oo$ that don't have the variables indexed by $I$.
\end{lemma}
\begin{proof} The discussion above provides the direction from left to right for $f = W$. To prove it for generic functions $f$ it suffices to note that $f_I^\oo$ and $f_J^\oo$ are independent from each other if $I \neq J$. That is, $D(f_I^\oo)$ and $D(f_J^\oo)$ do not share terms. The converse direction is proved by setting $f = W$ and setting $D = D'\mon{\ol{x}}{d}$ where $(D'\mon{\ol{x}}{d} )(f) = D'(\mon{\ol{x}}{d} f)$.
\end{proof}
This essentially turns this into a linear algebra problem. We just need to apply $D$ on generic functions $W_I^\oo$ that are missing the variables indexed by $I$, collect the different results and make their coefficients zero. That is, for every $f = W_I^\oo$ we look at the orbit
\begin{align}
	Orbit_G(f) = \{g(f)| g \in G\}
\end{align}
and collect coefficients in
\begin{align}
	0 = D(f) = \sum_{g\in G}R_g g(f) = \sum_{f' \in Orbit(f)} R_g' f'
\end{align}
where we set
\begin{align}
	R_g' = \sum_{g\in G, g(f) = f'} R_g = 0
\end{align}
Thus we obtain a linear system of $|Orbit_G(f)|$ equations in $|G|$ unknowns. Of course we do this for all $2^d - 1$ terms $f = W_I^\oo$ and hopefully obtain a system with a non-trivial solution. Here is a simple example
\begin{example} In $d = 2$ dimensions we have the generic functions $W(0,0) = const$ ,$W(0,x_2)$, $W(x_1,0)$ where at least one variable is missing. We can write them as $1, f(x_1), f(x_2)$. Now we consider the group generated by $g_k(x_k) = \ol{x}_k$
\begin{align}
	G = \{1, g_1, g_2, g_{12} = g_{21}\}
\end{align}
We'll write the operator $D$ as 
\begin{align}
	D = R - R_1 g_1 - R_2 g_2 + R_{12} g_{12}
\end{align}
We consider the different orbits. The one with the constant is trivial since the function doesn't change under the group action. Next we have
\begin{align}
	Orbit_G(f(x_1)) = \{f(x_1),f(\ol{x}_1),f(\ol{x}_1), f(\ol{x}_1)\} = \{f(x_1),f(\ol{x}_1)\}
\end{align}
and similarly
\begin{align}
	Orbit_G(f(x_2)) = \{f(x_2),f(\ol{x}_2)\}
\end{align}
This means we have the equations
\begin{align}
	R - R_1 &+ R_2 - R_{12} = 0\\
	R = R_2, &\quad R_1 = R_{12}\\
	R = R_1, &\quad R_2 = R_{12}
\end{align}
Folowing the implications provided from these operators we find $R= R_1 = R_2 = R_{12}$. The simplest choice is to set $R = 1$ and we find that every weighted orbit summation operator for this group $G$ must be trivially-weighted if it were to annihilate generic functions with missing variables.
\end{example}
\begin{definition} A weighted operator $D$ is trivial if all the weights are the same. Then we call the operator unweighted or regular.
\end{definition}
We have the following lemma
\begin{lemma} Let $G$ be a finite group with $d$ generators for a random walk model in $d$ dimensions with small steps and no interactions. Then any orbital operator is trivial.
\end{lemma}
We need to check this proof.
\begin{proof}It is enough to show that $D(W_k^0) = 0$ for all $k$ implies that the operator is trivial. To this end suppose that $D(f) = 0$ where $f = W_k^0$. Then we can partition the group $G$ as before into sets $H$ and $Hg_k$. We see that $hg_k (f) = h(f)$ so the orbit of $f$ contains at most $|G|/2$ elements. However these elements must all be distinct otherwise they would come from the same operations but the group elements in $H$ are all distinct. So $|Orbit(f) = |G|/2$. This means that we have the equations $R_h = R_{hg_k}$ for every $h \in H$. Once we go through every $k$ we have overlaps $R = R_1 =\cdots = R_d$ and $R_1 = R_{12} = R_{13}$ etc. This means that all the $R$'s are the same and the operator is trivial.
\end{proof}
This shows that in general there is no clever weighted summation operator that also satisfies the condition $D(\mon xd) \neq 0$. However, the idea will be fruitful in that we may find non-trivial partial solution in the interacting case.

%=======================================
\subsection{Walks with unconstrained steps}
%=======================================
We may go through the process again but this time we take
\begin{align}
	Q = \frac{W}{x_{[d]}^N K}
\end{align}
where $N_k$ is the maximum negative size of steps in the set $\mc{S}$ in the direction $x_k$ or in other words the maximum $N_k$ for which $S_k^{N_k}$ is non-zero. Then the operators that keep the walks in the quarter plane are
\begin{align}
	E_k = 1 - \sum_{i=1}^{N_k} \ext{\ol{x}_k^i}
\end{align}
Then we can conjugate as
\begin{align}
	E_k' = x_k^{N_k} E_k \ol{x}_k^{N_k} = 1 - \sum_{i=0}^{N_k-1}\ext{x_k^i} = 1 -\mc{E}_k
\end{align}
and so in the kernel equation we obtain
\begin{align}
	x_{[d]}^N = x_{[d]}^N E \ol{x}_{[d]}^N W = E' W = W +  \sum_{\emptyset \neq I\subseteq [d]} (-1)^{|I|} \mc{E}_I(W)
\end{align}
Here we have various powers $x_k^i$ along with terms with zeros in locations $I$ in some partial derivative of $W$. So if we seek an operator $D$ that does the same job as the orbit summation operator we need $D$ to annihilate all the terms with $x_k$'s missing in $W$ on the RHS. So them we'll have
\begin{align}
	D(x_{[d]}KQ) = D(W) = D\left(x_{[d]}^N-\sum_{\emptyset \neq I\subseteq [d]} (-1)^{|I|} \mc{E}_k(W)\right) = D(x_{[d]}^N)
\end{align}
In other words we want $D$ to behave (in some sense) similarly to the derivative
\begin{align}
\del_{[d]}^N = \del_1^{N_1}\cdots \del_d^{N_d}
\end{align}
This suggests some definitions.
\begin{definition}For a sequence $N = (N_1,\ldots, N_d)$ of positive integers we say that a model $\mc{M}$ is of type $N$ that the longest step in $\mc{S}$ in the negative $k$-th direction has length $N_k$ in that coordinate. In particular we may say that a model is of type $n$ if $N_k = n$ for all $k$.
\end{definition}
For example, full random walks with small steps are of type 1.
\begin{definition}For sequences $N = (N_1,\ldots, N_d)$ and $M = (M_1,\ldots, M_d)$ of positive integers we say that a model $\mc{M}$ is of type $(N,M)$ if it's of type $N$ and the longest step in $\mc{S}$ in the positive $k$-th direction has length $N_k$ in that coordinate. In particular we may say that a model is of type $(n,m)$ is of type $n$ and $M_k = m$ for all $k$.
\end{definition}
For example, full random walks with small steps are of type $1,1$. The notation is still a little awful. Anyway, it allows us to write things in the following form
\begin{theorem} A random walk model of type $N$ satisfies a functional equation
\begin{align}
	E'W = \fullx xdN
\end{align}
where $E' = \fullx xdN E \fullox xdN$ and $W  = \fullx xdN KQ$.
\end{theorem}

\begin{corollary} For a random walk model of type $N$ the generating function $Q$ for the non-interacting case satisfies the PDE
\begin{align}
	\fullx{\del}dN\left(\fullx xdN KQ\right) = 1
\end{align}
\end{corollary}
Here $\fullx{\del}dN$ is an example of operator $D$ that kills off all the unnecessary terms and doesn't kill off $\mon xd$. Talking about weighted operators here should be possible but it'll be more difficult. We may investigate type $(N,M)$ problems by finding a suitable model of this type whose orbital operator is nice enough. This will show that all type $(N,M)$ satisfy an orbital equation of some type, not just a PDE. We can consider the problem with $S = \sum_{k=1}^d \ol{x}_k^{n_k} + x_k^{n_k}$ of type $(N,N)$ for example whose group will be generated by the elements of groups $G_k$ with $2n_k$ elements each - each $g_{k,j,l} = \omega_{n_k}^jx_k^l$ where $\omega_{n_k} = \exp(2\pi i/n_k)$ and $l\in\{-1,1\}$. These elements commute and allow us to find an easy-enough orbital operator that does the job. Furthermore it does not kill off the $\fullx xdN$ term. This discussion suggest a deeper theory of which we may just give the first few definitions
\begin{definition} For a model $\mc{M}$ of type $N$ with pure operator $\mc{E} = 1- E'$ so that $W = \fullx xdN  +\mc{E}W$ let
\begin{align}
	X_{d,\mc{M}} = \{\text{Operators $D$}| D\mc{E} \equiv 0 \text{ on generic functions with $\leq d -1$ variables missing}\}
\end{align}
\end{definition}
This space is not empty since $\fullx{\del} dN \in X_{\mc{M}}$ as shown already. Also we have the orbital operator for type $N$ models.
\begin{definition} For any type $N$ model in $d$ dimensions define
\begin{align}
	X_d  = \{\text{Operators $D$}| D\mc{E} \equiv 0 \text{ on generic functions with $\leq d -1$ variables missing}\}
\end{align}
where $\mc{E}$ is for any model $\mc{M}$.
\end{definition}
We can define similar spaces $X_{d,\mc{M}}(F)$ and $X_d(F)$ where we want the extra property $D(Ff) = FD(f)$ for any function $f$. We can further define spaces $X^\ast$ where in addition to having that nice property on $K$ we have that $D(\fullx xdN) \neq 0$. This will help characterize all problems which have a useful orbit summation for example. And in the case where such summation does not exist we may look towards other useful operators that let us perform an analysis on the D-finiteness of $Q$.
%====================================
\section{Interactions case}
%====================================
It is natural to try extending the algebraic kernel method to the interacting case but we quickly see that the regular orbit summation operator is insufficient. This is where the weighted operator comes into play and allows us to say some things about certain cases of random walk models. The analysis for two-dimensional models is particularly nice because we end up having the same number of equations as unknowns in order to find the weights in the weighted operator. The determinants of the resulting matrices will tell us when we can have a partial solution using the weighted operator. Unfortunately we may not eliminate the $\Qf{0,0}$ variable from the equations. Thus the D-finiteness of the full function depends on the D-finiteness of the function $\Qf{0,0}$. We may be able to say something about the connections between the full interactions and the partial interactions via this paradigm. 

The functional equation in two dimensions is
\begin{align}
	KQ  - K_1 Q_1^0 -  K_2 Q_2^0 + K_{1,2}Q_{1,2}^{0,0} = \alpha
\end{align}
The group $G = G_S$ has two non-trivial generators $g_1$ and $g_2$ and $G$ has order 4, 6 or 8. We have the orbit summation operator $D_G$ which interprets the orbit summation method as an alternating sum of group actions.
%====================================
\subsection{Group of order 4}
%====================================
From a previous section we found that
\begin{align}
	D_G = 1-g_1-g_2+g_1g_2 = (1-g_1)(1-g_2) = \Delta_1\Delta_2
\end{align}
where we define $\Delta_i = 1- g_i$. In the unlabeled orbit summation method we had a direct application of $D_G$ on the equations to perform the required cancellations. Here we will attempt to do something similar. Since the cancellations don't happen automatically if we use $D_G$ directly, we'll instead obtain a linear system of equations by applying the group operations on the original functional equation. Then we'll use linear algebra to find a solution where our variables will be the ones we wish to eliminate $Q_i^0$ and their images under the group actions. If the associated matrix has non-zero determinant we have a unique solution and unfortunately we won't be able to eliminate variables. However there are some cases where the determinant will be zero. Because of symmetry we expect to be able eliminate the desired variables. The group operations commute so we have
\begin{align}
	g_i Q_i^0 = Q_i^0,\quad g_ig_j Q_i^0 = g_j Q_i^0
\end{align}
and therefore we have the variables that we wish to eliminate
\begin{align}
	E = \{Q_1^0,Q_2^0,g_2Q_1^0,g_1Q_2^0\}.
\end{align}
This is looking promising since we have four equations and four unknowns that we wish to eliminate. We re-write the original equation in a format centered on these variables
\begin{align}
	 K_1 Q_1^0 +  K_2 Q_2^0 = KQ + K_{1,2}Q_{1,2}^{0,0} -\alpha = \beta
\end{align}
For brevity let $g_{12} = g_1g_2$
\begin{align}
	 K_1 Q_1^0 +  K_2 Q_2^0 &= \beta\\
	 g_1K_1 Q_1^0 + g_1K_2 g_1Q_2^0 &= g_1\beta = \beta_1\\
	 g_2K_1 g_2Q_2^0 + g_2K_2 Q_2^0 &= g_2\beta = \beta_2\\
	 g_{12}K_1 g_2Q_1^0 + g_{12}K_2 g_1Q_2^0 &= g_{12}\beta = \beta_{12}\\
\end{align}
Using the displayed order of variables in $E$ as our order for the system we can write the associated matrix $M$ as
\begin{align}
	M_4 = \begin{bmatrix}
	K_1 & K_2 & 0 & 0 \\
	g_1 K_1 & 0 & 0 & g_1 K_2\\
	0 & g_2 K_2 & g_2 K_1 & 0\\
	0 & 0 & g_{12}K_1 & g_{12} K_2
	\end{bmatrix}
\end{align}
Its determinant is
\begin{align}
	\det M_4 &= K_1 g_{12}K_1 g_1K_2g_2K_2 -g_1K_1g_2K_1K_2g_{12}K_2\\
	& = K_1 g_{12}K_1 g_1K_2g_2K_2  -g_1(g_1K_1g_{12}K_1g_1K_2g_{2}K_2)\\
	&=\Delta_1  (K_1 g_{12}K_1 g_1K_2g_2K_2)\\
	&=\Delta_2 (g_1 K_1 g_2 K_1K_2 g_{12} K_2 )
\end{align}
where the last form comes from doing the process in the same way but with $g_2$ instead.

%====================================
\subsubsection{When is $\det M_4 = 0$?}
%====================================

  Now our task is to find where this determinant is zero. Clearly the expression withing each $\Delta$ must be in the kernel of the operator. To that end let's find when that happens. For a Laurent polynomial $f$ in the variables $x_i$ and $t$ we have
\begin{align}
	\Delta_k f = \sum_n f_n (1-g_k)x_k^n = \sum_n (f_n - f_{-n}\mu_k^{-n})x_i^n
\end{align}
where $\mu_k = S_k^{-1}/S_k^1$. This implies that
\begin{align}
	f_{n} &= f_{-n} \mu_k^{-n}\\
	f_{n} (S_k^1)^n &= f_{-n}(S_k^{-1})^n
\end{align}
for all $n \in \ze$. It's not clear if this will help us. However the models that exhibit groups of order 4 have at least one symmetry across a line - either vertical or horizontal symmetry. WLOG let the model have vertical symmetry. That is 
\begin{align}
	g_1(x_1) = x_1, \quad g_1(S_2^i) = S_2^i,\quad g_1(S_1^i) = S_1^{-i}
\end{align}
This allows us to compute some things. For example
\begin{align}
	g_1 K_2 &= c_2 - t\ol{x}_2 g_1S_2^{-1} = c_2 - t\ol{x}_2 S_2^{-1} = K_2
	g_1g_2K_2 = g_2 g_1 K_2 = g_2
\end{align}
Therefore $\Delta_1$ fixes both $K_2$ and $g_2K_2$.
We see that the determinant is
\begin{align}
	\det M_4 = \Delta_1  (K_1K_2 g_{12}K_1 g_2K_2) = K_2 g_2K_2 \Delta_1  (K_1 g_{12}K_1)
\end{align}
This is zero if and only if $\Delta_1  (K_1 g_{12}K_1) = 0$ so we have
\begin{align}
	0&=\Delta_1  (K_1 g_{12}K_1) \\
	&= K_1g_{12}K_1 - g_1 K_1 g_2 K_2 \\
	&= (c_1-t\ol{x_1}S_1^{-1})(c_1-tx_1g_2S_1^{-1})-(c_1-tx_1S_1^{-1})(c_1-t\ol{x_1}g_2S_1^{-1})\\
	&= c_1 t \ol{x}_1 (g_2 S_2^{-1} - S_2^{-1}) - c_1 t x_1 (g_2 S_2^{-1} - S_2^{-1})\\
	&= c_1 t (\ol{x}_1-x_1)(g_2 S_1^{-1}-S_1^{-1})\\
	& = -c_1 t(\ol{x}_1-x_1)\Delta_2(S_1^{-1})
\end{align}
This is true if and only if
\begin{align}
	g_2 S_1^{-1} = S_1^{-1}
\end{align}
We have
\begin{align}
	g_2 (\Sf{-1,-1}\ol{x}_2+\Sf{-1,0}+\Sf{-1,1}x_2) = \Sf{-1,-1}x_2 \ol{\mu}_2+\Sf{-1,0}+\Sf{-1,1}\ol{x}_2\mu_2
\end{align}
and
\begin{align}
	\Sf{-1,-1}x_2 \ol{\mu}_2+\Sf{-1,1}\ol{x}_2\mu_2 = \Sf{-1,-1}\ol{x}_2+\Sf{-1,1}x_2
\end{align}
which can be re-written succinctly as
\begin{align}
	\Sf{-1,-1} = \mu_2 S_{-1,1}
\end{align}
If one of the $S$'s is zero so is the other and therefore by symmetry the model doesn't have diagonal steps. Since we have two non-trivial generators this implies that both $x_2$ and $\ol{x}_2$ are present in $S$, so the model trivially exhibits a symmetry about the $x_1$ axis. If the $S$'s are non-zero this implies $\mu_2$ is a constant and therefore $S_2^{-1} = S_2^1$. Again the model exhibits symmetry about the $x_1$ axis. This leads to a theorem of sorts. The determinant of $M_4$ is zero if and only if the model is highly-symmetric. These are the only cases where we have hope that the orbit summation method will extend to the labeled cases. Note that in the unlabeled case $c_1,c_2 = 0$ so the analysis yields a zero determinant in any group 4 case, confirming that we can use the orbit summation method.

%====================================
\subsubsection{Analogue of the Orbit Summation Method}
%====================================

We are now ready to attempt using the orbit summation method in the highly symmetric cases. There are 8 such models but only 5 of them have dimension 2. Furthermore 2 of the models are reflections of each other so we are really left with 4 interesting cases - diagonal walks, full walks, full walks minus two horizontal steps. If we are able to eliminate all the variables then we'll be left with a linear relation
\begin{align}
	\gamma \beta- \gamma_1 \beta_1-\gamma_2 \beta_2 + \gamma_{12}\beta_{12} = 0
\end{align}
where the choice of signs will make things neater later. Let $\bullet = g_1$ and $\ast = g_2$. First we have $K_2^\bullet = K_2$ and $K_1^\ast = K_1$. Then we let
\begin{align}
	x = Q_1^0, y= Q_2^0, u= Q_1^{0,\ast}, v = Q_2^{0,\bullet}
\end{align}
and
\begin{align}
	a = K_1, b = K_1^\bullet,c=K_2,d = K_2^\ast
\end{align}
Then the equations from the system can be described as
\begin{align}
	\beta &= a x + c y,\quad \beta_1 = b x + c v\\
	\beta_2 &= d y + a u,\quad \beta_{12} = b u + d v
\end{align}
Therefore the above equation needs to satisfy
\begin{align}
	0 &= (ax+cy)\gamma - (bx+cv)\gamma_1 - (dy+au)\gamma_2 + (bu +dv) \gamma_{12}\\
	   &= (a\gamma-b\gamma_1)x + (c\gamma-d\gamma_2)y + (b\gamma_{12}-a\gamma_2)u + (d\gamma_{12}-c\gamma_1)v
\end{align}
and in order to eliminate the variables $x,y,u,v$ we need to make all their coefficients zero. Therefore
\begin{align}
	a\gamma = b\gamma_1,\quad c\gamma = d\gamma_2, \quad b\gamma_{12} = a \gamma_2,\quad d\gamma_{12} = c\gamma_1
\end{align}
If we set $\gamma = 1$ then $\gamma_1 = a/b$, $\gamma_2 = c/d$ and $\gamma_{12} = (ac)/(bd)$. So we can multiply out by $bd$ to find the relation
\begin{align}
	bd\beta-ad\beta_1-cb \beta_2+ac \beta_{12} = 0
\end{align}
where we have successfully eliminated $Q_1^0, Q_2^0, Q_1^{0,\ast}$, and $  Q_2^{0,\bullet}$. Or written in terms of the $K$'s we have
\begin{align}
	K_1^{\bullet}K_2^{\ast}\beta^{}-K_1^{}K_2^{\ast}\beta^{\bullet}-K_1^{\bullet}K_2^{}\beta^{\ast}+K_1^{}K_2^{}\beta^{\bullet\ast} = 0
\end{align}
Written in operator notation this equation becomes much nicer
\begin{align}
	(1-\bullet)(1-\ast)(K_1^\bullet K_2^\ast \beta) = \Delta(K_1^\bullet K_2^\ast \beta) = 0
\end{align}
Let $K_0 = K_1^\bullet K_2^\ast$. We can now write this equation in another way
\begin{align}
	\Delta(K_0\beta) &= \Delta(K_0(KQ + K_{1,2}Q_{1,2}^{0,0} - \alpha)) \\
	&= \Delta(K_0 KQ) +\Delta(K_0 K_{1,2}Q_{1,2}^{0,0}) - \Delta(K_0\alpha)\\
	&= K\Delta(K_0 Q) + \Delta(K_0 K_{1,2})Q_{1,2}^{0,0} - \Delta(K_0)\alpha
\end{align}
where the last line follows from $\Delta$ fixing $K$ and functions without variables $x_i$. All that remains is to calculate the various quantities.
\begin{align}
	\Delta g_iK_i = - \Delta K_i = - (c_i-tx_iS_i^{-1}-c_i + t\ol{x}_i S_i^{-1}) = -t(\ol{x}_i-x_i)S_i^{-1}
\end{align}
Therefore
\begin{align}
	\Delta K_0 = \Delta_1 (K_1^\bullet) \Delta_2 (K_2^\ast) = t^2(\ol{x}_1-x_1)(\ol{x}_2-x_2)S_1^{-1}S_2^{-1}
\end{align}
Next
\begin{align}
	\Delta(K_0 K_{1,2}) = \Delta_1 K_1^\bullet \Delta_2 K_2^\ast K_{1,2} = \Delta_1 K_1^\bullet (K_2^\ast K_{1,2}- K_2 K_{1,2}^\ast)
\end{align}
We calculate the quantity in parentheses as
\begin{align}
	&K_2^\ast K_{1,2}- K_2 K_{1,2}^\ast =\\
	=& (c_2-tx_2S_2^{-1})(c_1c_2-t\ol{x_1}\ol{x_2}S_{1,2}^{-1,-1})-(c_2-t\ol{x_2}S_2^{-1})(c_1c_2-t\ol{x_1}x_2S_{1,2}^{-1,-1})\\
	=&c_1c_2tS_2^{-1}(\ol{x_2}-x_2) - c_2t\ol{x_2}S_{1,2}^{-1,-1}(\ol{x_2}-x_2)\\
	=&c_2t(\ol{x_2}-x_2)(c_1S_2^{-1}-\ol{x_1}S_{1,2}^{-1,-1})
\end{align}
and thus
\begin{align}
	\Delta(K_0 K_{1,2})  &= c_2t(\ol{x_2}-x_2) \Delta_1 K_1^\bullet (c_1S_2^{-1}-\ol{x_1}S_{1,2}^{-1,-1})\\
	&=c_2t(\ol{x_2}-x_2) \Delta_1 K_1^\bullet (c_1S_2^{-1}-\ol{x_1}S_{1,2}^{-1,-1})\\
\end{align}
To compute this note that
\begin{align}
\Delta_1 K_1^\bullet c_1 S_2^{-1} = c_1 S_2^{-1}\Delta_1 K_1^\bullet = -c_1t S_2^{-1}(\ol{x}_1-x_1)
\end{align}
and
\begin{align}
	\Delta_1 K_1^\bullet \ol{x_1}S_{1,2}^{-1,-1} = S_{1,2}^{-1,-1}\Delta_1 K_1^\bullet \ol{x_1} = c_1 (\ol{x}_1-x_1)S_{1,2}^{-1,-1}
\end{align}
Finally
\begin{align}
	\Delta(K_0 K_{1,2}) = c_1c_2t(\ol{x}_1-x_1)(\ol{x}_2-x_2)(t S_1^{-1} S_2^{-2}-S_{1,2}^{-1,-1}).
\end{align}
Combining terms and manipulating things a bit leads us to the equation
\begin{align}
	\Delta(K_0 Q) = \alpha t(\ol{x}_1-x_1)(\ol{x}_2-x_2)(tS_1^{-1}S_2^{-1} - b_1b_2(tS_1^{-1}S_2^{-1}-\Sf{-1,-1})\Qf{0,0})
\end{align}
or
\begin{align}
	\Delta(K_0 Q) = \alpha t(\ol{x}_1-x_1)(\ol{x}_2-x_2)(tS_1^{-1}S_2^{-1}(1 - b_1b_2\Qf{0,0}) + \Sf{-1,-1}\Qf{0,0})
\end{align}

%====================================
\subsection{Similarities between labeled and unlabeled versions}
%====================================

We start with the observation that $\Delta \equiv D_G$ for the case of highly-symmetric random walk models
\begin{align}
	D_G(K_1^\bullet K_2^\ast(KQ + K_{1,2}Q_{1,2}^{0,0} - \alpha)) = 0
\end{align}
If we move down the the unlabelled case (assuming we can do with through this equation) we have
\begin{align}
	K_1^\bullet \mapsto - t x_1 S_1^{-1}\\
	K_2^\ast \mapsto - t x_2 S_2^{-1}\\
	K_{1,2} \mapsto - t \ol{x}_1\ol{x}_2 S_{1,2}^{-1,-1}
\end{align}
and the equation becomes
\begin{align}
	t^2S_1^{-1}S_2^{-1}D_G(x_1x_2(KQ -t \ol{x}_1\ol{x}_2S_{1,2}^{-1,-1}Q_{1,2}^{0,0}-1)
\end{align}
since $\Delta$ fixes $S_i^{-1}$. The $\Qf{0,0}$ quantity disappears after the action of $D_G$ and we arrive at
\begin{align}
	D_G'(KQ-1) = 0
\end{align}
which was the orbit sum equation.
%====================================
\subsection{Group of order 6}
%====================================

Here we have some interactions as well but the problem is harder. We still have that $g_i Q_i^0 = Q_i^0$. We need to see what kind of expressions we get from the $Q_i^0$. Here we tabulate the sequences of operations $g_1$ and $g_2$ and which parts of those sequences survive application on $Q_1^0,Q_2^0, K_1$ and $K_2$. For example
\begin{align}
g_{2121}Q_1^0 = g_{212}Q_1^0 = g_{121} Q_1^0 = g_{12}Q_1^0
\end{align}
so the table has $12$ in the appropriate rown/column.

\begin{center}
    \begin{tabular}{| l | l | l | l | l | l |}
    \hline
    Sequence & Reverse & Q1 & Q2 & K1 & K2 \\ \hline
    1 & 21212 & n/a & 1& 1& 1 \\ \hline
    21 &  1212& 2& 21 & 21& 21  \\ \hline
    121 & 212 & 12& 21& 121& 212 \\ \hline
    2121 & 12 &  12 & 1 & 12 & 12 \\ \hline
    12121 & 2 &  2& n/a& 2& 2 \\
    \hline
    \end{tabular}
\end{center}

As the talbe shows we really have only 6 variables. Let
\begin{align}
	u_1 &= Q_1^0,\quad u_2 = g_2 Q_1^0,\quad u_3 = g_{12}Q_1^0\\
	v_1 &= Q_2^0,\quad v_2 = g_1 Q_2^0,\quad v_3 = g_{21}Q_2^0\\
\end{align}
Also there are 10 different $K_1$'s and $K_2$'s which we label sequentially as $A_1,\ldots,A_6$ and $B_1,\ldots,B_6$ respectively. We have the associated $ 6 \times 6$ matrix
\begin{align}
	M_6 = \begin{bmatrix}
	A_1& & & B_1 & & \\
	A_2& & & & B_2& \\
	& A_3& & & & B_3\\
	& & A_4 & & & B_4\\
	& & A_5& & B_5& \\
	&A_6 & &B_6 & & \\
	\end{bmatrix}
\end{align}
with determinant
\begin{align}
	\det M_6 = A_2A_4A_6 B_1 B_3 B_5 - A_1 A_3 A_5 B_2 B_4 B_6
\end{align}
What is nice about these coefficients is that $A_1 = g_1 A_2$ and so on and thus we can write
\begin{align}
	\det M_6 = \Delta_1 (A_2A_4A_6 B_1 B_3 B_5) = \Delta_1(g_1K_1 g_{121}K_1g_{12}K_1 K_2 g_{21}K_2 g_{12}K_2)
\end{align}
We can easily check the relationships $A$'s and $B$'s in the two group cases. For the models with the first group
\begin{itemize}
\item $S = x_1\ol{x}_2+\ol{x}_1+x_2$ : $A_1 = A_6, A_2 = A_5, A_3 = A_4$ and $B_1 = B_4, B_2 = B_3, B_5 = B_6$.
\item $S = \ol{x}_1x_2 + x_1 + \ol{x}_2$:  $A_1 = A_4, A_2 = A_3, A_5 = A_6$ and $B_1 = B_2, B_3 = B_6, B_4 = B_5$
\item $S = x_1+\ol{x}_1+x_2+\ol{x}_2+ x_1\ol{x}_2 + \ol{x}_1 x_2$: We have no non-trivial relations.
\end{itemize}
For the models with the second group we have
\begin{itemize}
\item $S = x_1 x_2+\ol{x}_1+\ol{x}_2$ :  $A_1 = A_6, A_2 = A_5, A_3 = A_4$ and $B_1 = B_2, B_3 = B_6, B_4 = B_5$
\item $S = \ol{x}_1\ol{x}_2 + x_1 + x_2$: $A_1 = A_4, A_2 = A_3, A_5 = A_6$ and $B_1 = B_4, B_2 = B_3, B_5 = B_6$.
\item $S = x_1+\ol{x}_1+x_2+\ol{x}_2+ x_1x_2 + \ol{x}_1\ol{x}_2$: We have no non-trivial relations.
\end{itemize}
For the cases with relations it is clear that $\det M_6 = 0$. Using a computer we can easily check that the determinant is not zero in the cases without nontrivial relations. It's worth mentioning that so far all groups that have had determinant zero have been crystallographic but as we can see not all crystallographic groups end up having determinant zero. 

We seek the same type of relation on $\beta$ and its orbit as before so we'll use $D_G$ as inspiration. We have the following
\begin{align}
0=&\gamma\beta-\gamma_{1}\beta_{1}+\gamma_{21}\beta_{21}-\gamma_{121}\beta_{121}+\gamma_{12}\beta_{12}-\gamma_{2}\beta_{2}\\
	=&\gamma (A_1 u_1 + B_1 v_1) - \gamma_{2} (A_6 u_2 + B_6 v_1) - \gamma_{1} A_2 u_1 + B_2 v_2) \\&+  \gamma_{12} (A_5 u_3 + B_5 v_2) + \gamma_{21}(A_3 u_2 + B_3 v_3) - \gamma_{121} (A_4 u_3 + B_4 v_3).
\end{align}
Expanding out and collecting terms leads to the following equations
\begin{align}
	 A_1 \gamma&= A_2 \gamma_{1},   \quad A_3 \gamma_{21} = A_6 \gamma_{2}, \quad A_5 \gamma_{12} = A_4 \gamma_{121}\\
	 B_1 \gamma &= B_6\gamma_{2},  \quad  B_5 \gamma_{12}= B_2 \gamma_{1},  \quad B_3 \gamma_{21} = B_4 \gamma_{121}
\end{align}
This system of equations doesn't have a nontrivial symbolic solution so we need to use the relations here. There is a rational solution in each case with relations where $\gamma$ is a free variable. We set $\gamma$ to be the least common multiple of the denominators. In the six cases we've performed the calculations and checked that indeed $K_0 = \gamma$ in each case.
\begin{itemize}
\item $S = x_1\ol{x}_2+\ol{x}_1+x_2$ : $\quad\gamma = K_0 =  A_3 A_5 B_6 = g_{21}(K_1) g_{12}(K_1) g_{2}(K_2)$
\item $S = \ol{x}_1x_2 + x_1 + \ol{x}_2$:   $\quad\gamma = K_0 =  A_2 A_5 B_6 = g_{1}(K_1) g_{12}(K_1) g_{2}(K_2)$
\item $S = x_1 x_2+\ol{x}_1+\ol{x}_2$ :  $\quad\gamma = K_0 =  A_2 A_3 B_5 B_6 = g_{1}(K_1) g_{21}(K_1) g_{12}(K_2) g_{2}(K_2)$
\item $S = \ol{x}_1\ol{x}_2 + x_1 + x_2$:  $\quad\gamma = K_0 =  A_2 B_6 = g_{1}(K_1) g_{2}(K_2)$
\end{itemize}
Here we've picked the respresentations of $A$'s and $B$' with the least number of group operations $g_i$ to reach the resulting expression. We also tried finding a possible $K_0$ for the cases with no relations but we didn't find any among products of distinct $A$'s and $B$'s.
%====================================
\subsection{Group of order 8}
%====================================
The analysis for groups of order 8 is quite similar so we shall omit it. Here are the results. 
The two models are
\begin{align}
	S = x_1+\ol{x}_1+x_1\ol{x}_2+\ol{x}_1x_2\quad\quad S  = x_1+\ol{x}_1+x_1x_2+\ol{x}_1\ol{x}_2
\end{align}
The relations are the same in both models
\begin{align}
	A_1 &= A_6, A_2 = A_5, A_3 = A_4, A_7 = A_8\\
	B_1 &= B_4, B_2 = B_3, B_4 = B_8, B_6 = B_7
\end{align}
In both cases
\begin{align}
	\gamma = K_0 =  A_2 A_3 B_5 B_6 = g_1(K_1) g_{21}(K_1) g_{212}(K_2) g_{1212}(K_2)
\end{align}
Here's we've also written them in their shortest forms.

%====================================
\subsection{Infinite Groups}
%====================================

A natural question is whether we can extend this analysis to handle infinite groups, at least in a purely formal sense. Here the orbits will be infinitely long and while there will be relations between the $g Q_i^0$ variables from one orbit iteration to the next, there will be no point where we can stop that we'll have as many equations as variables - there will always be more variables. Thus we'll need to take an infinite summation of the equations. It's quite possible that any candidates for $K_0$ will be infinite products. Since the $K_i$'s have $t^0$ coefficients, such products will not converge. However, there may be hope if there are enough relations between the $g K_i$'s so that we may have a finite product of $g K_i$'s for $K_0$. The unlabelled orbit method doesn't need a $K_0$ to get rid of all terms $Q_i^0$ and $\Qf{0,0}$. One would hope that something like this would work for the infinite case.

\endinput